\documentclass[12pt]{article}
\title{Response Memo}

\begin{document}
\maketitle
 	 \paragraph{Clarifying the main thesis.-} I followed reviewer D's recommendation to be more explicit about the fact that aesthetic experience does not necessarily result in ``knowledge,'' since this knowledge can be gained by multiple sources besides direct experience, and in fact routinely is, especially by members of culturally and economically privileged status groups.  This resulted in a major re-organization of the argument offered in the front-end of the paper, which I think has made the basic thesis of the paper sharper and more succinct.  Now it is clear that it is precisely the existence of residual stratification by status net of consumption that is the most significant finding here. I thank the reviewer for pointing to this flaw in the original argument and offering guidance in this respect. 
      \paragraph{Interactive versus additive specifications.-}Reviewer D suggested a comparison between an ``additive model'' (in which consumption frequency enters as an additional control variable in predicting talk) and an interactive model in which the effects of education and other socio-demographic predictors are allowed to vary across levels of consumption.  I build the conceptual justification for these different ways of parsing the data into the (now re-rewritten) theoretical discussion. The results of including interaction effects between selected socio-demographic predictors and frequency of consumption in predicting culture-talk are now shown in Table 2 .  As shown by the paucity of statistically significant results, the additive model is certainly much more applicable to the data than the interactive model.
      \paragraph{Important matters.-} Reviewers A, B and D recommended excising the section on ``important matters.''  I have followed the reviewers' recommendation in this regard. Instead, the end of the paper is now dedicated with exploring the wider implications of the results for inequalities in social connectivity connected to status.
      \paragraph{Simmel and socialibility.-} Reviewer D suggested that the link to Simmelian sociability in the previous version was not effective, since the link to stratification is not germane to this idea. In this version of the paper, I have kept the reference to culture-talk and sociability but now I explicitly that thinking of it as just Simmelian sociability is not enough.  I then use this as a platform to more clearly show how we can link the sociability idea to the culture and stratification literature.

\end{document}•
